\documentclass[11pt]{article}

\usepackage[T1]{fontenc}
\usepackage{lmodern}
\usepackage[margin=1in]{geometry}
\usepackage{amsmath,amssymb,amsthm,mathtools}
\usepackage{microtype}
\usepackage[hidelinks]{hyperref}
\usepackage{enumitem}

\setlength{\parindent}{0pt}
\setlength{\parskip}{0.35em}
\setlength{\emergencystretch}{3em}
\setlist{nosep}

\theoremstyle{definition}
\newtheorem{definition}{Definition}[section]
\newtheorem{exercise}{Exercise}[section]

\newcommand{\R}{\mathbb R}
\newcommand{\C}{\mathbb C}
\newcommand{\Q}{\mathbb Q}
\newcommand{\Z}{\mathbb Z}
\newcommand{\N}{\mathbb N}
\newcommand{\T}{\mathbb T}
\newcommand{\id}{\operatorname{id}}
\newcommand{\im}{\operatorname{im}}
\newcommand{\Tr}{\operatorname{Tr}}
\newcommand{\End}{\operatorname{End}}
\newcommand{\Hom}{\operatorname{Hom}}
\newcommand{\supp}{\operatorname{supp}}
\newcommand{\spec}{\operatorname{spec}}
\newcommand{\coker}{\operatorname{coker}}
\newcommand{\rank}{\operatorname{rank}}
\newcommand{\vol}{\operatorname{vol}}
\newcommand{\ad}{\operatorname{ad}}
\newcommand{\Ad}{\operatorname{Ad}}

\hypersetup{
  pdftitle={A Physical Introduction to Higher Mathematics},
  pdfauthor={Paul Gustafson}
}

\title{A Physical Introduction to Higher Mathematics}
\author{Paul Gustafson}
\date{}

\begin{document}
\maketitle
\tableofcontents
\clearpage

\part{Waves}

\section{Linear Structure and Spectra}

\begin{definition}[Vector spaces and linear maps]
Let $k$ be a field. A \emph{$k$-vector space} is an abelian group $V$ with an
action $k\times V\to V$ satisfying distributivity, associativity, and
$1v=v$. A map $T:V\to W$ is \emph{linear} if
$T(av+bw)=aT(v)+bT(w)$.
\end{definition}

\begin{definition}[Subspaces, kernels, images, and quotients]
A \emph{subspace} of $V$ is a subset closed under linear combinations. For a
linear map $T:V\to W$, define
\[
\ker T=\{v:T(v)=0\},\qquad \im T=\{T(v):v\in V\}.
\]
For a subspace $U\subseteq V$, the \emph{quotient} $V/U$ is the vector space
of cosets $v+U$. The \emph{direct sum} $V\oplus W$ is $V\times W$ with
componentwise operations.
\end{definition}

\begin{exercise}[Universal linear constructions]
Prove that $\ker T$ and $\im T$ are subspaces. State and prove the universal
properties of $V/U$ and $V\oplus W$. Deduce
$V/\ker T\cong\im T$ and the three isomorphism theorems.
\end{exercise}

\begin{definition}[Duals and pairings]
The \emph{dual} of $V$ is $V^*=\Hom_k(V,k)$. A \emph{pairing} of $V$ and
$W$ is a bilinear map $V\times W\to k$. It is \emph{nondegenerate} if each
nonzero argument pairs nontrivially with some argument in the other space.
\end{definition}

\begin{definition}[Tensor products]
A \emph{tensor product} of $V$ and $W$ is a vector space $V\otimes W$ with a
bilinear map $(v,w)\mapsto v\otimes w$ such that every bilinear map
$b:V\times W\to X$ factors uniquely through a linear map
$\widetilde b:V\otimes W\to X$. Tensor powers and multilinear maps are
defined by iteration.
\end{definition}

\begin{exercise}[Duality and multilinearity]
For finite-dimensional spaces, construct the natural maps
\[
V^*\otimes W\longrightarrow\Hom(V,W),\qquad
V^*\otimes W^*\longrightarrow(V\otimes W)^*.
\]
Prove that they are isomorphisms, identify the evaluation pairing, and prove
the tensor--hom adjunction
$\Hom(U\otimes V,W)\cong\Hom(U,\Hom(V,W))$.
\end{exercise}

\begin{definition}[Inner products and orthogonality]
An \emph{inner product} on a real vector space is a positive-definite
symmetric bilinear form. On a complex vector space it is a positive-definite
sesquilinear form, conjugate-linear in the first variable. Write
$\|v\|=\sqrt{\langle v,v\rangle}$. Vectors are \emph{orthogonal} when their
inner product is zero, and
$U^\perp=\{v:\langle u,v\rangle=0\text{ for every }u\in U\}$.
\end{definition}

\begin{definition}[Projections and adjoints]
A \emph{projection} is an endomorphism $P$ with $P^2=P$. An \emph{orthogonal
projection} is a projection with $\ker P=(\im P)^\perp$. The \emph{adjoint}
of $T:V\to W$ is a map $T^*:W\to V$ satisfying
$\langle Tv,w\rangle=\langle v,T^*w\rangle$.
\end{definition}

\begin{exercise}[Orthogonal decomposition and best approximation]
Let $V$ be finite-dimensional and $U\subseteq V$. Prove
$V=U\oplus U^\perp$, construct the orthogonal projection $P_U$, and prove
that $P_Uv$ is the unique element of $U$ minimizing $\|v-u\|$. Prove that
$P$ is an orthogonal projection if and only if $P=P^*=P^2$.
\end{exercise}

\begin{definition}[Isometries and spectral classes]
A linear map is \emph{orthogonal} over $\R$, or \emph{unitary} over $\C$, if
it preserves the inner product. An endomorphism $A$ is \emph{self-adjoint} if
$A=A^*$ and \emph{normal} if $AA^*=A^*A$. A scalar $\lambda$ is an
\emph{eigenvalue} if $\ker(A-\lambda\id)\ne0$; a subspace $U$ is
\emph{invariant} if $A(U)\subseteq U$.
\end{definition}

\begin{definition}[Spectral decomposition]
A \emph{spectral decomposition} of an endomorphism $A$ is an expression
\[
A=\sum_{\lambda\in\spec(A)}\lambda P_\lambda,
\qquad P_\lambda P_\mu=\delta_{\lambda\mu}P_\lambda,
\qquad \sum_\lambda P_\lambda=\id,
\]
with orthogonal projections $P_\lambda$.
\end{definition}

\begin{exercise}[Finite-dimensional spectral theorem]
Prove that adjoints exist uniquely. Prove that a complex endomorphism is
unitarily diagonalizable if and only if it is normal. Deduce spectral
decompositions for self-adjoint, unitary, and normal operators, including the
reality of self-adjoint spectra and the unit modulus of unitary spectra.
\end{exercise}

\begin{exercise}[Diagonalization and invariant subspaces]
For an endomorphism with split characteristic polynomial, prove that it is
diagonalizable if and only if the space is the direct sum of its eigenspaces.
Show that normality removes generalized eigenvectors. Prove that commuting
normal endomorphisms admit a common orthogonal spectral decomposition.
\end{exercise}

\section{Oscillation and Fourier Analysis}

\begin{definition}[Linear oscillator]
A \emph{linear oscillator} on a finite-dimensional real inner-product space
$V$ is an equation
\[
M\ddot q+Kq=0,
\]
where $M$ is positive definite and $K$ is self-adjoint and positive
semidefinite. A \emph{normal mode} is a nonzero $v\in V$ for which
$q(t)=a(t)v$ solves the equation for some nonconstant $a$.
\end{definition}

\begin{exercise}[The harmonic oscillator]
Solve $m\ddot q+kq=0$ for $m,k>0$. Prove conservation of
$E=\frac12m\dot q^2+\frac12kq^2$, determine the period, and deduce from the
measured period of a small-amplitude mass--spring system the ratio $k/m$.
\end{exercise}

\begin{exercise}[Coupled oscillators and normal modes]
Transform $M\ddot q+Kq=0$ by $M^{1/2}$. Apply the spectral theorem to prove
the existence of an $M$-orthogonal basis of normal modes and derive their
frequencies. For a fixed-end chain of $N$ equal masses joined by equal
springs, compute the frequencies and mode shapes and identify the modes seen
in a Chladni-type standing-wave measurement.
\end{exercise}

\begin{definition}[Vibrating string]
For length $L$, tension $\tau$, and linear density $\rho$, a \emph{vibrating
string} is a function $u:[0,L]\times\R\to\R$ satisfying fixed-end conditions
$u(0,t)=u(L,t)=0$ and
\[
\rho\,\partial_t^2u=\tau\,\partial_x^2u.
\]
\end{definition}

\begin{exercise}[String modes and measured pitch]
Derive the vibrating-string equation from transverse force balance in the
small-slope limit. Determine all separated normal modes and their
frequencies. Prove that the fundamental frequency is
$f_1=(2L)^{-1}\sqrt{\tau/\rho}$ and compare the predicted ratios
$f_n/f_1=n$ with the harmonics of a stretched string.
\end{exercise}

\begin{definition}[Orthogonal expansions]
Let $H$ be an inner-product space. For an orthogonal family $(e_n)$, the
\emph{orthogonal coefficients} of $f$ are
$\widehat f(n)=\langle e_n,f\rangle/\|e_n\|^2$. A \emph{Fourier series} on
$\T=\R/(2\pi\Z)$ is the formal expansion
$f\sim\sum_{n\in\Z}\widehat f(n)e^{inx}$, where
$\widehat f(n)=(2\pi)^{-1}\int_0^{2\pi}f(x)e^{-inx}\,dx$.
\end{definition}

\begin{exercise}[Bessel, completeness, and Parseval]
Prove Bessel's inequality for every finite orthogonal family. Prove that an
orthonormal family $(e_n)$ is complete precisely when
\[
\|f\|^2=\sum_n|\langle e_n,f\rangle|^2
\]
for every $f$. Prove completeness of $(e^{inx})_{n\in\Z}$ in
$L^2(\T)$ and obtain Parseval's identity and $L^2$ convergence of Fourier
series.
\end{exercise}

\begin{exercise}[Spectral evolution on an interval]
Expand suitable fixed-end initial data in string modes and derive the unique
solution of the wave equation. Expand suitable initial temperature data in
the same spatial eigenfunctions and derive the fixed-end solution of
$\partial_tu=\kappa\partial_x^2u$. Prove conservation of wave energy and
monotone decay of heat energy directly from the spectral coefficients.
\end{exercise}

\begin{definition}[Fourier transform and convolution]
For $f\in L^1(\R^n)$, its \emph{Fourier transform} is
\[
\mathcal Ff(\xi)=(2\pi)^{-n/2}\int_{\R^n}f(x)e^{-ix\cdot\xi}\,dx.
\]
For $f,g\in L^1(\R^n)$, their \emph{convolution} is
$(f*g)(x)=\int_{\R^n}f(y)g(x-y)\,dy$.
\end{definition}

\begin{exercise}[Fourier inversion, convolution, and Plancherel]
For Schwartz functions, prove Fourier inversion and
$\mathcal F(f*g)=(2\pi)^{n/2}(\mathcal Ff)(\mathcal Fg)$. Prove that
$\mathcal F$ preserves the $L^2$ inner product and extends uniquely to a
unitary operator on $L^2(\R^n)$. Deduce the Plancherel theorem.
\end{exercise}

\begin{exercise}[Wave and heat propagators]
Use the Fourier transform to solve the Cauchy problems for
\[
\partial_t^2u=c^2\Delta u,\qquad \partial_tu=\kappa\Delta u
\]
on $\R^n$. Identify their spectral-mode evolution, derive the Gaussian heat
kernel and its convolution law, and prove finite propagation speed for the
one-dimensional wave equation.
\end{exercise}

\begin{exercise}[Interference and Fourier optics]
For coherent scalar waves with complex amplitudes $a_1,a_2$, derive
$|a_1+a_2|^2$ and the fringe spacing in Young's double-slit geometry. In the
Fraunhofer regime, derive that the far-field amplitude is the Fourier
transform of the aperture. Compute the single-slit diffraction intensity,
locate its zeros, and compare the predicted angular spacing with measured
optical diffraction.
\end{exercise}

\part{Symmetry}

\section{Groups and Representations}

\begin{definition}[Groups, homomorphisms, and quotients]
A \emph{group} is a set $G$ with an associative multiplication, an identity,
and inverses. A \emph{homomorphism} $\phi:G\to H$ preserves multiplication.
A subgroup $N\leq G$ is \emph{normal} if $gNg^{-1}=N$ for every $g\in G$;
the \emph{quotient group} $G/N$ consists of cosets with induced
multiplication.
\end{definition}

\begin{exercise}[The group isomorphism theorems]
Prove that kernels are normal and that a normal subgroup is the kernel of the
quotient homomorphism. Prove the three isomorphism theorems for groups and
the universal property of $G/N$.
\end{exercise}

\begin{definition}[Actions, orbits, and stabilizers]
A \emph{left action} of $G$ on $X$ is a homomorphism
$G\to\operatorname{Bij}(X)$. The \emph{orbit} and \emph{stabilizer} of
$x\in X$ are $Gx=\{gx:g\in G\}$ and $G_x=\{g:gx=x\}$.
\end{definition}

\begin{exercise}[Orbit--stabilizer and class equation]
Construct a natural bijection $G/G_x\to Gx$. For finite $G$, deduce the
orbit--stabilizer formula and the class equation from the conjugation action.
Use these results to prove that every group of prime order is cyclic.
\end{exercise}

\begin{definition}[Representations and intertwiners]
A \emph{representation} of $G$ on $V$ is a homomorphism
$\rho:G\to\operatorname{GL}(V)$. An \emph{intertwiner}
$T:(V,\rho)\to(W,\sigma)$ satisfies $T\rho(g)=\sigma(g)T$. A subspace is
\emph{invariant} if it is preserved by every $\rho(g)$; a nonzero
representation is \emph{irreducible} if it has no proper nonzero invariant
subspace. Direct sums and tensor products carry the actions
$\rho\oplus\sigma$ and $g\mapsto\rho(g)\otimes\sigma(g)$.
\end{definition}

\begin{exercise}[Maschke's theorem]
Let $G$ be finite and let the characteristic of $k$ not divide $|G|$.
Average an arbitrary projection over $G$ to construct an invariant
complement to every invariant subspace. Deduce that every finite-dimensional
representation is a direct sum of irreducible representations.
\end{exercise}

\begin{exercise}[Schur's lemma]
Prove that a nonzero intertwiner between irreducible complex
finite-dimensional representations is an isomorphism and that every
endomorphism of an irreducible representation is scalar. Deduce uniqueness
of irreducible multiplicities in a decomposition.
\end{exercise}

\begin{definition}[Characters]
The \emph{character} of a finite-dimensional representation $V$ is the class
function $\chi_V(g)=\Tr(\rho_V(g))$. For class functions on a finite group,
define
\[
\langle f,h\rangle_G=\frac1{|G|}\sum_{g\in G}\overline{f(g)}h(g).
\]
\end{definition}

\begin{exercise}[Character orthogonality and decomposition]
Use averaging of linear maps and Schur's lemma to prove orthonormality of the
irreducible characters. Prove that they form an orthonormal basis of the
class functions and that
$\langle\chi_U,\chi_V\rangle_G=\dim\Hom_G(U,V)$. Deduce the multiplicity
formula and the character table constraints.
\end{exercise}

\begin{exercise}[Symmetry projections]
For an irreducible character $\chi$ of a finite group, prove that
\[
P_\chi=\frac{\dim\chi}{|G|}\sum_{g\in G}\overline{\chi(g)}\rho(g)
\]
is the projection onto the $\chi$-isotypic subspace. Apply these projections
to the displacement space of a molecule whose equilibrium configuration is
preserved by $G$.
\end{exercise}

\section{Continuous Symmetry and Angular Momentum}

\begin{definition}[Lie groups and Lie algebras]
A \emph{Lie group} is a smooth manifold $G$ whose multiplication and
inversion are smooth. Its \emph{Lie algebra} is
$\mathfrak g=T_eG$ with bracket induced by left-invariant vector fields. A
\emph{one-parameter subgroup} is a smooth homomorphism $\R\to G$. The
\emph{exponential map} $\exp:\mathfrak g\to G$ sends $X$ to the value at
$1$ of the one-parameter subgroup with initial velocity $X$.
\end{definition}

\begin{exercise}[One-parameter subgroups and the exponential map]
Prove existence and uniqueness of the one-parameter subgroup generated by
$X\in\mathfrak g$. For a closed matrix group, prove that it is
$t\mapsto e^{tX}$, identify the commutator bracket, and compute the first
nontrivial terms of the Baker--Campbell--Hausdorff formula.
\end{exercise}

\begin{definition}[Infinitesimal representations]
For a smooth representation $\rho:G\to\operatorname{GL}(V)$, its
\emph{infinitesimal representation} is
\[
d\rho:\mathfrak g\to\End(V),\qquad
d\rho(X)=\left.\frac{d}{dt}\right|_{0}\rho(\exp tX).
\]
\end{definition}

\begin{exercise}[Differentiating symmetry]
Prove that $d\rho$ preserves brackets and that
$\rho(\exp X)=\exp(d\rho(X))$. Prove that a connected Lie-group
representation is determined by its infinitesimal representation whenever
the latter integrates uniquely.
\end{exercise}

\begin{definition}[Rotations and spin]
Define
\[
\begin{aligned}
SO(3)&=\{R\in\End(\R^3):R^*R=\id,\ \det R=1\},\\
SU(2)&=\{U\in\End(\C^2):U^*U=\id,\ \det U=1\}.
\end{aligned}
\]
An \emph{angular-momentum representation} is a unitary representation of
$SU(2)$; with Planck constant $\hbar$, its infinitesimal generators
$J_1,J_2,J_3$ are normalized by
$[J_i,J_j]=i\hbar\varepsilon_{ijk}J_k$.
\end{definition}

\begin{exercise}[$SU(2)$ double-covers $SO(3)$]
Let $SU(2)$ act by conjugation on the real space of traceless Hermitian
$2\times2$ operators. Prove that the resulting homomorphism
$SU(2)\to SO(3)$ is surjective with kernel $\{\pm\id\}$. Deduce the effect of
a $2\pi$ rotation in integer- and half-integer-spin representations.
\end{exercise}

\begin{exercise}[Irreducible representations of $SU(2)$]
Using raising and lowering operators, prove that every finite-dimensional
irreducible complex representation is indexed by
$j\in\{0,\frac12,1,\ldots\}$, has dimension $2j+1$, and has
$J_3$-eigenvalues $m\hbar$ for $m=-j,-j+1,\ldots,j$. Compute the spectrum of
$J^2=J_1^2+J_2^2+J_3^2$.
\end{exercise}

\begin{exercise}[Clebsch--Gordan decomposition]
Prove
\[
V_{j_1}\otimes V_{j_2}\cong
\bigoplus_{j=|j_1-j_2|}^{j_1+j_2}V_j.
\]
Construct coupled angular-momentum eigenstates from highest-weight vectors
and compute the singlet and triplet decomposition of
$V_{1/2}\otimes V_{1/2}$.
\end{exercise}

\begin{definition}[Quantum symmetries]
For a Hilbert space $H$ and self-adjoint Hamiltonian $\mathcal H$, a
\emph{unitary symmetry} is a unitary $U$ satisfying
$U\mathcal H=\mathcal HU$. A \emph{quantum number} is an eigenvalue of a
self-adjoint operator commuting with $\mathcal H$. An operator $A$ has
\emph{symmetry type} $W$ when its components define an intertwiner from $W$
to $\End(H)$ under conjugation.
\end{definition}

\begin{exercise}[Symmetry, degeneracy, and selection rules]
Prove that every eigenspace of $\mathcal H$ is symmetry-invariant and
decomposes into irreducibles. Deduce symmetry-forced degeneracies. Prove that
a transition matrix element $\langle f,Ai\rangle$ can be nonzero only if the
trivial representation occurs in
$V_f^*\otimes W\otimes V_i$, and derive the angular-momentum rules for an
electric-dipole operator.
\end{exercise}

\begin{exercise}[Atomic spectra]
For the Coulomb Hamiltonian, use rotational symmetry and separation into
irreducible $SO(3)$-types to obtain the quantum numbers $\ell,m$ and their
degeneracies. With the radial spectrum
$E_n=-m_e e^4/(2(4\pi\varepsilon_0)^2\hbar^2n^2)$ as derived input, compute
the Balmer wavelengths and compare them with the observed hydrogen lines.
Apply the electric-dipole rule to identify forbidden transitions.
\end{exercise}

\begin{exercise}[Molecular symmetry]
For the equilibrium geometry of $H_2O$, determine the point-group action on
nuclear displacements, remove translations and rotations, and decompose the
vibrational representation into irreducibles. Determine infrared and Raman
activity from the transformation types of the dipole and polarizability.
Compare the predicted number and symmetry of lines with the observed three
fundamental vibrational bands.
\end{exercise}

\part{Fields}

\section{Manifolds, Forms, and Metrics}

\begin{definition}[Topological and smooth manifolds]
A \emph{topological manifold} of dimension $n$ is a Hausdorff,
second-countable space locally homeomorphic to $\R^n$. A \emph{smooth
manifold} is a topological manifold with a maximal atlas whose transition
maps are smooth. A map of smooth manifolds is \emph{smooth} when its chart
representatives are smooth.
\end{definition}

\begin{definition}[Tangent and cotangent spaces]
The \emph{tangent space} $T_pM$ is the vector space of derivations
$v:C^\infty(M)\to\R$ at $p$, satisfying
$v(fg)=f(p)v(g)+g(p)v(f)$. The \emph{cotangent space} is
$T_p^*M=(T_pM)^*$. The differential of $F:M\to N$ is
$dF_p(v)(f)=v(f\circ F)$.
\end{definition}

\begin{exercise}[Intrinsic tangent vectors]
Prove that derivations at $p$, equivalence classes of smooth curves through
$p$, and the usual coordinate tangent vectors are naturally isomorphic.
Prove the chain rule $d(G\circ F)_p=dG_{F(p)}\circ dF_p$ and compute tangent
spaces of regular level sets.
\end{exercise}

\begin{definition}[Vector fields and flows]
A \emph{vector field} is a smooth section $X:M\to TM$. A \emph{flow} of
$X$ is a smooth local action $\Phi$ of $\R$ on $M$ satisfying
$\left.\frac d{dt}\right|_0\Phi_t(p)=X_p$. The \emph{Lie bracket} is the
vector field $[X,Y](f)=X(Yf)-Y(Xf)$.
\end{definition}

\begin{exercise}[Flows and brackets]
Prove local existence and uniqueness of flows. Prove that $[X,Y]=0$ if and
only if their local flows commute, and that the Lie derivative
$\mathcal L_XY=[X,Y]$ is natural under diffeomorphisms.
\end{exercise}

\begin{definition}[Exterior algebra and differential forms]
The \emph{exterior algebra} of a vector space $V$ is
$\Lambda V=T(V)/\langle v\otimes v:v\in V\rangle$. A \emph{differential
$k$-form} on $M$ is a smooth section of $\Lambda^kT^*M$. The \emph{pullback}
of $\omega\in\Omega^k(N)$ by $F:M\to N$ is
\[
(F^*\omega)_p(v_1,\ldots,v_k)
=\omega_{F(p)}(dF_pv_1,\ldots,dF_pv_k).
\]
\end{definition}

\begin{definition}[Exterior derivative]
The \emph{exterior derivative} is the unique family of maps
$d:\Omega^k(M)\to\Omega^{k+1}(M)$ satisfying $df(X)=Xf$, the graded Leibniz
rule, $d^2=0$, and naturality $F^*d=dF^*$.
\end{definition}

\begin{exercise}[Cartan calculus]
Construct $d$ locally and prove its uniqueness. For contraction $\iota_X$,
prove Cartan's identity
$\mathcal L_X=d\iota_X+\iota_Xd$ and the naturality of $d$, wedge products,
and pullbacks.
\end{exercise}

\begin{definition}[Integration of forms]
An \emph{orientation} of an $n$-manifold is a continuous choice of component
of nonzero elements in $\Lambda^nT_p^*M$. The integral
$\int_M\omega$ of a compactly supported top form is defined using
orientation-preserving charts and a partition of unity.
\end{definition}

\begin{exercise}[Stokes' theorem]
Prove that the integral is independent of charts and partition of unity.
For an oriented manifold with boundary, prove
\[
\int_Md\omega=\int_{\partial M}\omega.
\]
Deduce the fundamental theorem of calculus, Green's theorem, the divergence
theorem, and the classical curl theorem.
\end{exercise}

\begin{definition}[Tensor fields and metrics]
A \emph{$(r,s)$-tensor field} is a smooth section of
$TM^{\otimes r}\otimes T^*M^{\otimes s}$. A \emph{Riemannian metric} is a
smooth positive-definite inner product on each $T_pM$. A \emph{Lorentzian
metric} is a smooth nondegenerate symmetric form of signature
$(-,+,\ldots,+)$. The length of a curve in a Riemannian manifold is
$\int\sqrt{g(\dot\gamma,\dot\gamma)}\,dt$; the proper time of a timelike
curve is $c^{-1}\int\sqrt{-g(\dot\gamma,\dot\gamma)}\,dt$. A metric and an
orientation define a volume form $\vol_g$.
\end{definition}

\begin{exercise}[Metric volume and proper time]
Construct $\vol_g$ intrinsically and derive its coordinate expression.
Prove invariance of length and proper time under reparametrization. In
Minkowski spacetime, maximize proper time among inertial paths with fixed
endpoints and derive special-relativistic time dilation.
\end{exercise}

\section{Connections, Curvature, and Field Equations}

\begin{definition}[Connections and geodesics]
A \emph{connection} on $TM$ is a map
$\nabla:\Gamma(TM)\times\Gamma(TM)\to\Gamma(TM)$, linear over smooth
functions in the first variable and satisfying
$\nabla_X(fY)=X(f)Y+f\nabla_XY$. Its torsion is
$T(X,Y)=\nabla_XY-\nabla_YX-[X,Y]$. A \emph{geodesic} satisfies
$\nabla_{\dot\gamma}\dot\gamma=0$. A connection is \emph{metric compatible}
if $Xg(Y,Z)=g(\nabla_XY,Z)+g(Y,\nabla_XZ)$.
\end{definition}

\begin{exercise}[Levi--Civita connection]
Prove that every Riemannian or Lorentzian metric has a unique torsion-free,
metric-compatible connection. Derive the Koszul formula and prove that its
geodesics are precisely the critical curves of the energy functional with
fixed endpoints.
\end{exercise}

\begin{definition}[Covariant derivative and parallel transport]
For a curve $\gamma$, the connection induces a \emph{covariant derivative}
$D/dt$ on vector fields along $\gamma$. A field $V$ is \emph{parallel} when
$DV/dt=0$; the resulting linear isomorphism between endpoint tangent spaces
is \emph{parallel transport}.
\end{definition}

\begin{exercise}[Parallel transport and holonomy]
Prove existence, uniqueness, composition, and inverse properties of parallel
transport. Show that Levi--Civita transport preserves the metric. Compute the
rotation obtained by transporting a tangent vector around a latitude on the
unit sphere and relate it to the enclosed curvature.
\end{exercise}

\begin{definition}[Curvature]
The \emph{curvature} of a connection is
\[
R(X,Y)Z=\nabla_X\nabla_YZ-\nabla_Y\nabla_XZ-\nabla_{[X,Y]}Z.
\]
For a metric connection, the \emph{Riemann tensor} is
$\operatorname{Riem}(X,Y,Z,W)=g(R(X,Y)Z,W)$. Its contraction is the
\emph{Ricci tensor} $\operatorname{Ric}$; the contraction of Ricci is the
\emph{scalar curvature} $S$.
\end{definition}

\begin{exercise}[Curvature identities and geodesic deviation]
Prove tensoriality and the algebraic symmetries of the Riemann tensor and the
first and second Bianchi identities. For a variation through geodesics with
separation field $J$, prove
\[
\frac{D^2J}{dt^2}+R(J,\dot\gamma)\dot\gamma=0.
\]
Deduce the leading relative acceleration of nearby freely falling bodies.
\end{exercise}

\begin{definition}[Gauge fields]
Let $G$ be a Lie group and $\pi:P\to M$ a principal $G$-bundle. A
\emph{gauge connection} is an equivariant $\mathfrak g$-valued one-form
$A$ on $P$ reproducing fundamental vertical fields. Its \emph{curvature} is
$F_A=dA+\frac12[A\wedge A]$. A \emph{gauge transformation} is a
$G$-equivariant bundle automorphism covering $\id_M$.
\end{definition}

\begin{exercise}[Gauge covariance and Bianchi identity]
Derive the transformation laws for a local gauge potential and its
curvature. Prove $d_AF_A=0$ and gauge invariance of holonomy up to conjugacy.
For $G=U(1)$, recover $A\mapsto A+d\chi$ and $F=dA$.
\end{exercise}

\begin{definition}[Electromagnetic field]
On an oriented Lorentzian four-manifold, an \emph{electromagnetic field
strength} is a two-form $F$. Given a current one-form $J$ and constants
$\mu_0,c$, \emph{Maxwell's equations} are
\[
dF=0,\qquad d{*F}=\mu_0{*J}.
\]
Locally, a \emph{gauge potential} is a one-form $A$ with $F=dA$.
\end{definition}

\begin{exercise}[Maxwell theory and electromagnetic waves]
Derive charge conservation from Maxwell's equations. On Minkowski
spacetime, recover the electric and magnetic vector equations and the Lorentz
force from $F$. In vacuum, impose Lorenz gauge and derive wave equations with
speed $c$. Prove transverse polarization and compare the predicted relation
$c=(\mu_0\varepsilon_0)^{-1/2}$ with the measured speed of light.
\end{exercise}

\begin{definition}[Einstein equation]
The \emph{Einstein tensor} of a Lorentzian metric is
$G=\operatorname{Ric}-\frac12Sg$. A \emph{stress--energy tensor} is a
symmetric two-tensor $T$ whose contraction with an observer's velocity gives
the observer's energy--momentum current. The \emph{Einstein equation} with
cosmological constant $\Lambda$ is
\[
G+\Lambda g=\frac{8\pi G_N}{c^4}T.
\]
\end{definition}

\begin{exercise}[Conservation and the Newtonian limit]
Use the contracted Bianchi identity to derive $\nabla^aT_{ab}=0$. Linearize
$g=\eta+h$ for a static weak field and recover Poisson's equation and the
Newtonian inverse-square law with the stated coupling constant.
\end{exercise}

\begin{exercise}[Gravitational radiation]
Linearize the vacuum Einstein equation about Minkowski spacetime, impose
transverse--traceless gauge, and obtain two tensor polarizations satisfying a
wave equation. Apply geodesic deviation to a ring of freely falling test
masses and derive the strain measured by an orthogonal-arm interferometer.
\end{exercise}

\begin{exercise}[Gravitational-wave detection]
For a quasi-circular binary, derive the leading quadrupole waveform and the
chirp relation between frequency, frequency derivative, and chirp mass.
Using public values for GW150914, infer the chirp mass and order of magnitude
of the strain. Compare these predictions with the measured coincident LIGO
strain, limiting the conclusion to agreement of the waveform with general
relativity in the observed regime.
\end{exercise}

\part{States}

\section{Hamiltonian and Statistical States}

\begin{definition}[Configuration and phase spaces]
A \emph{configuration space} is a smooth manifold $Q$. Its \emph{phase
space} is the cotangent bundle $T^*Q$. The tautological one-form on $T^*Q$ is
$\theta_\alpha(v)=\alpha(d\pi(v))$, and its canonical symplectic form is
$\omega=-d\theta$.
\end{definition}

\begin{definition}[Symplectic dynamics]
A \emph{symplectic form} is a closed, nondegenerate two-form $\omega$. The
\emph{Hamiltonian vector field} of $H\in C^\infty(M)$ is defined by
$\iota_{X_H}\omega=dH$. Its flow is the \emph{Hamiltonian flow}. The
\emph{Poisson bracket} is $\{f,g\}=\omega(X_f,X_g)$.
\end{definition}

\begin{exercise}[Hamilton's equations and Poisson algebra]
Prove that the canonical form on $T^*Q$ is symplectic. In canonical
coordinates, derive Hamilton's equations. Prove antisymmetry, the Leibniz
rule, and the Jacobi identity for the Poisson bracket, and prove
$\dot f=\{H,f\}+\partial_tf$ along Hamiltonian flow.
\end{exercise}

\begin{definition}[Symplectic actions and momentum maps]
An action of a Lie group $G$ on $(M,\omega)$ is \emph{symplectic} if it
preserves $\omega$. A \emph{momentum map} is an equivariant map
$J:M\to\mathfrak g^*$ such that
$d\langle J,\xi\rangle=\iota_{\xi_M}\omega$ for every
$\xi\in\mathfrak g$.
\end{definition}

\begin{exercise}[Noether's theorem]
Prove that if a Hamiltonian is invariant under a symplectic $G$-action with
momentum map $J$, then every component of $J$ is conserved. Compute the
momentum maps for translations and rotations of a particle and recover
linear and angular momentum conservation.
\end{exercise}

\begin{exercise}[Liouville's theorem]
On a $2n$-dimensional symplectic manifold, prove that Hamiltonian flow
preserves $\omega$, the phase-space volume $\omega^n/n!$, and every density
depending only on $H$. Deduce that Hamiltonian evolution cannot compress
phase volume.
\end{exercise}

\begin{definition}[Measure and integration]
A \emph{$\sigma$-algebra} $\Sigma$ on $X$ is closed under complements and
countable unions. A \emph{measure} is a countably additive map
$\mu:\Sigma\to[0,\infty]$ with $\mu(\varnothing)=0$. A function is
\emph{measurable} if inverse images of Borel sets are measurable. Its
integral is obtained from nonnegative simple functions by monotone limits and
then by positive and negative parts. For $1\leq p<\infty$,
$L^p(X,\mu)$ consists of measurable functions modulo equality almost
everywhere with $\int|f|^p\,d\mu<\infty$; $L^\infty$ uses essential bounds.
\end{definition}

\begin{exercise}[Convergence and $L^p$ spaces]
Prove monotone and dominated convergence and Fatou's lemma. Prove H\"older's
and Minkowski's inequalities and completeness of $L^p$. Prove that $L^2$ is
a Hilbert space and identify the dual of $L^p$ for $1<p<\infty$ under the
standard $\sigma$-finiteness hypotheses.
\end{exercise}

\begin{definition}[Probability]
A \emph{probability space} is a measure space $(\Omega,\mathcal F,\mathbb P)$
with $\mathbb P(\Omega)=1$. A \emph{random variable} is a measurable map.
Its \emph{expectation} is its integral. The \emph{conditional expectation}
$\mathbb E[X\mid\mathcal G]$ is the $\mathcal G$-measurable integrable random
variable satisfying
$\int_G\mathbb E[X\mid\mathcal G]d\mathbb P=\int_GX,d\mathbb P$ for all
$G\in\mathcal G$. Families of $\sigma$-algebras are \emph{independent} when
the probability of every finite intersection is the product of its
probabilities.
\end{definition}

\begin{exercise}[Expectation, conditioning, and independence]
Prove existence and almost-everywhere uniqueness of conditional expectation,
the tower property, and conditional Jensen inequality. Characterize
independence by factorization of expectations and derive variance addition
for independent square-integrable random variables.
\end{exercise}

\begin{definition}[Entropy and Gibbs ensembles]
For a discrete probability vector $p$, its \emph{entropy} is
$S(p)=-k_B\sum_ip_i\log p_i$. For a phase-space probability density $\rho$
relative to Liouville measure $\mu_L$, its \emph{statistical entropy} is
$S(\rho)=-k_B\int\rho\log\rho,d\mu_L$. The \emph{canonical Gibbs ensemble}
at inverse temperature $\beta$ is
$d\mathbb P_\beta=Z(\beta)^{-1}e^{-\beta H}d\mu_L$.
\end{definition}

\begin{definition}[Measure-preserving dynamics and ergodicity]
A measurable transformation $T$ is \emph{measure preserving} if
$\mu(T^{-1}A)=\mu(A)$. It is \emph{ergodic} if every invariant measurable
set has measure zero or full measure. An \emph{equilibrium state} is an
invariant probability measure for the dynamics.
\end{definition}

\begin{exercise}[Entropy and equilibrium]
Maximize entropy under normalization and fixed mean energy to derive the
Gibbs ensemble. Prove
$\langle H\rangle=-\partial_\beta\log Z$ and
$\operatorname{Var}(H)=\partial_\beta^2\log Z$. Prove invariance of Gibbs
measure under Hamiltonian flow and formulate the consequence of the ergodic
theorem for equality of time and ensemble averages.
\end{exercise}

\begin{exercise}[Ideal-gas thermodynamics]
Compute the classical canonical partition function of $N$ indistinguishable
noninteracting particles in a box, including the phase-space normalization
$h^{3N}N!$. Derive the ideal-gas law, internal energy, and entropy, and
compare the pressure law with dilute-gas measurements within the classical
regime.
\end{exercise}

\section{Hilbert-Space and Algebraic States}

\begin{definition}[Hilbert spaces and bounded operators]
A \emph{Hilbert space} is a complete inner-product space. A linear map
$T:H\to K$ is \emph{bounded} if
$\|T\|=\sup_{\|x\|\leq1}\|Tx\|<\infty$. Its adjoint is the unique bounded
map satisfying $\langle Tx,y\rangle=\langle x,T^*y\rangle$. Projections,
self-adjoint operators, and unitary operators satisfy respectively
$P=P^*=P^2$, $A=A^*$, and $U^*U=UU^*=\id$.
\end{definition}

\begin{definition}[Projection-valued measures]
A \emph{projection-valued measure} on a measurable space $X$ is a map
$E$ from measurable sets to orthogonal projections on $H$ with
$E(X)=\id$ and countable additivity in the strong operator topology. For a
bounded measurable $f$, the operator $\int_Xf,dE$ is defined by uniform
approximation from simple functions.
\end{definition}

\begin{exercise}[Spectral theorem]
Prove that every bounded self-adjoint operator $A$ has a unique
projection-valued measure on $\spec(A)\subset\R$ such that
$A=\int\lambda,dE(\lambda)$. Extend the statement to bounded normal
operators, derive continuous functional calculus, and characterize unitary
and projection spectra.
\end{exercise}

\begin{definition}[Quantum states and observables]
A \emph{density operator} is a positive trace-class operator $\rho$ with
$\Tr\rho=1$. An \emph{observable} is a self-adjoint operator $A$ with
spectral measure $E_A$. The \emph{Born probability} of an event $B$ is
$\Tr(\rho E_A(B))$. A closed-system \emph{unitary dynamics} is
$\rho\mapsto U_t\rho U_t^*$ for a strongly continuous unitary group
$U_t$.
\end{definition}

\begin{exercise}[Born rule and unitary dynamics]
Prove that the Born rule defines a probability measure and that its expected
value is $\Tr(\rho A)$ when defined. Derive the von Neumann equation from
$U_t=e^{-itH/\hbar}$. Prove conservation of expected energy and derive the
Heisenberg equation for observables.
\end{exercise}

\begin{definition}[Composite and reduced states]
The state space of a composite system is $H_A\otimes H_B$. A state is
\emph{separable} if it lies in the closed convex hull of product density
operators and \emph{entangled} otherwise. The \emph{reduced state}
$\rho_A=\Tr_B\rho$ is determined by
$\Tr(\rho_AA)=\Tr(\rho(A\otimes\id))$ for all bounded $A$.
\end{definition}

\begin{exercise}[Schmidt decomposition and entanglement]
Prove the Schmidt decomposition for pure vectors in a tensor product of
finite-dimensional Hilbert spaces. Prove that a pure state is separable
exactly when its reduced state has rank one, and that the two reduced states
have the same nonzero spectrum. Compute the reduced states and entanglement
entropy of a Bell state.
\end{exercise}

\begin{definition}[$C^*$-algebras and states]
A \emph{Banach algebra} is a complete normed algebra with
$\|ab\|\leq\|a\|\|b\|$. An \emph{involution} satisfies
$(ab)^*=b^*a^*$ and $(a^*)^*=a$. A \emph{$C^*$-algebra} is an involutive
Banach algebra satisfying $\|a^*a\|=\|a\|^2$. An element is \emph{positive}
if it is $b^*b$ for some $b$. A \emph{state} is a positive linear functional
$\varphi$ with $\|\varphi\|=\varphi(1)=1$.
\end{definition}

\begin{exercise}[Commutative Gelfand duality]
For compact Hausdorff $X$, prove that $C(X)$ is a commutative unital
$C^*$-algebra. Prove that every commutative unital $C^*$-algebra $A$ is
isometrically $*$-isomorphic to $C(\operatorname{Spec}A)$ by the Gelfand
transform, and identify states on $C(X)$ with probability measures.
\end{exercise}

\begin{definition}[Representations and GNS construction]
A \emph{representation} of a $C^*$-algebra $A$ is a $*$-homomorphism
$\pi:A\to B(H)$. For a state $\varphi$, define
$N_\varphi=\{a:\varphi(a^*a)=0\}$ and
$\langle[a],[b]\rangle=\varphi(a^*b)$ on $A/N_\varphi$.
\end{definition}

\begin{exercise}[GNS theorem]
Complete $A/N_\varphi$ and prove that left multiplication induces a cyclic
representation $(\pi_\varphi,H_\varphi,\Omega_\varphi)$ satisfying
$\varphi(a)=\langle\Omega_\varphi,\pi_\varphi(a)\Omega_\varphi\rangle$.
Prove uniqueness up to a cyclic unitary equivalence.
\end{exercise}

\begin{definition}[Operator topologies and von Neumann algebras]
The \emph{weak operator topology} is generated by
$T\mapsto\langle x,Ty\rangle$; the \emph{strong operator topology} is
generated by $T\mapsto\|Tx\|$. For $S\subseteq B(H)$, its \emph{commutant}
is $S'=\{T:TS=ST\text{ for all }S\in S\}$. A \emph{von Neumann algebra} is
a unital $*$-subalgebra $M\subseteq B(H)$ with $M=M''$.
\end{definition}

\begin{exercise}[Bicommutant theorem]
Prove that for a unital $*$-subalgebra $A\subseteq B(H)$, the strong and weak
operator closures of $A$ equal $A''$. Deduce that von Neumann algebras have
sufficient projections to approximate their self-adjoint elements
spectrally.
\end{exercise}

\begin{definition}[Traces, factors, and noncommutative probability]
A \emph{trace} on a von Neumann algebra is a normal positive functional
$\tau$ satisfying $\tau(ab)=\tau(ba)$. A \emph{factor} is a von Neumann
algebra with center $\C1$. A \emph{noncommutative probability space} is a
pair $(M,\varphi)$ of a von Neumann algebra and a normal state; its random
variables are operators affiliated with $M$.
\end{definition}

\begin{exercise}[Projection dimension and factors]
For a finite factor with normalized trace $\tau$, prove additivity and unitary
invariance of $\tau$ on projections. Relate the possible projection values
to matrix factors and to diffuse finite factors, and formulate independence
of commuting subalgebras by factorization of the state.
\end{exercise}

\section{Completely Bounded Structure and Correlations}

\begin{definition}[Positive and completely positive maps]
A linear map $\Phi:A\to B$ between $C^*$-algebras is \emph{positive} if it
preserves positive elements and \emph{completely positive} if every
$\id_{M_n}\otimes\Phi$ is positive. A \emph{quantum channel} is a completely
positive trace-preserving map on trace-class operators, equivalently a
unital completely positive map on observables in the adjoint picture.
\end{definition}

\begin{exercise}[Stinespring and Kraus representations]
Prove that a unital completely positive map $\Phi:A\to B(H)$ has a
representation
$\Phi(a)=V^*\pi(a)V$. In finite dimensions, derive a Kraus representation
$\Phi(a)=\sum_iK_i^*aK_i$ and its normalization. Prove that every channel is
unitary evolution on a larger system followed by discarding an environment.
\end{exercise}

\begin{definition}[Measurements]
A \emph{positive operator-valued measure} is a countably additive map $E$
from measurable events to positive operators with $E(X)=1$. In state $\rho$,
its outcome probability is $\Tr(\rho E(B))$.
\end{definition}

\begin{exercise}[POVMs and dilation]
Prove that projection-valued measures are POVMs. Prove Naimark's dilation
theorem in finite dimensions and derive the measurement channel associated
to a POVM. Show that state discrimination can require a nonprojective POVM.
\end{exercise}

\begin{definition}[Operator spaces]
An \emph{operator space} is a closed subspace $E\subseteq B(H)$ with matrix
norms inherited from $M_n(E)\subseteq B(H^{\oplus n})$. A map
$u:E\to F$ is \emph{completely bounded} if
$\|u\|_{cb}=\sup_n\|\id_{M_n}\otimes u\|<\infty$. The matrix order of an
operator system is the family of positive cones $M_n(E)_+$.
\end{definition}

\begin{exercise}[Ruan's theorem]
Prove that concrete operator-space norms satisfy
\[
\|x\oplus y\|=\max(\|x\|,\|y\|),\qquad
\|axb\|\leq\|a\|\|x\|\|b\|.
\]
Prove conversely that every vector space with matrix norms satisfying these
axioms admits a complete isometry into $B(H)$.
\end{exercise}

\begin{exercise}[Arveson--Wittstock extension]
Prove in finite dimensions that a completely positive map from an operator
system extends completely positively to the containing $C^*$-algebra.
Deduce, by the $2\times2$ operator-system construction, that completely
bounded maps admit norm-preserving extensions.
\end{exercise}

\begin{definition}[Operator-space tensor products]
The \emph{minimal operator-space tensor norm} on $E\otimes F$ is inherited
from concrete inclusions in $B(H\otimes K)$. The \emph{Haagerup norm} is
\[
\|z\|_h=\inf\{\|[x_1\ \cdots\ x_r]\|\,
\|[y_1\ \cdots\ y_r]^T\|:z=\sum_ix_i\otimes y_i\}.
\]
A bilinear form $u:E\times F\to\C$ is \emph{jointly completely bounded}
when its associated matrix amplifications are uniformly bounded.
\end{definition}

\begin{exercise}[Tensor norms linearize bilinear maps]
Prove the universal property of the projective Banach-space tensor norm.
Prove that the Haagerup tensor product linearizes completely bounded
factorizations and compare the minimal, Haagerup, and projective norms on
finite-dimensional matrix spaces.
\end{exercise}

\begin{exercise}[Grothendieck inequalities]
State and prove the finite-dimensional real Grothendieck inequality for
bilinear forms on $\ell_\infty^m\times\ell_\infty^n$. Formulate the
noncommutative and operator-space Grothendieck inequalities and deduce their
dimension-independent comparison of bounded and completely bounded
correlations.
\end{exercise}

\begin{definition}[Bell functionals and XOR games]
A \emph{Bell functional} is an array $M=(M_{xy}^{ab})$ evaluated on
conditional distributions by
$\langle M,p\rangle=\sum_{a,b,x,y}M_{xy}^{ab}p(a,b\mid x,y)$. An \emph{XOR
game} is specified by real coefficients $M_{xy}$ and scores correlations
$\sum_{x,y}M_{xy}\mathbb E[a_xb_y]$ for signs $a_x,b_y$.
\end{definition}

\begin{exercise}[Classical and quantum correlation bounds]
Prove that the classical and quantum XOR values are
\[
\sup_{a_x,b_y=\pm1}\left|\sum M_{xy}a_xb_y\right|,
\qquad
\sup_{\|u_x\|,\|v_y\|\leq1}
\left|\sum M_{xy}\langle u_x,v_y\rangle\right|.
\]
Apply Grothendieck's inequality to bound their ratio. For CHSH, compute the
classical bound $2$ and quantum bound $2\sqrt2$.
\end{exercise}

\begin{exercise}[Measured nonclassical correlations]
For a polarization-entangled two-photon state, choose four polarization
angles attaining the CHSH value $2\sqrt2$ and derive the predicted
coincidence correlations. Analyze a loophole-free Bell-test data table with
its stated uncertainties and determine whether the classical bound is
excluded, restricting the conclusion to local hidden-variable models
satisfying the experimental assumptions.
\end{exercise}

\part{Topology}

\section{Homotopy, Cohomology, and Bundles}

\begin{definition}[Homotopy and fundamental group]
A \emph{homotopy} from $f:X\to Y$ to $g:X\to Y$ is a continuous map
$H:X\times[0,1]\to Y$ with endpoints $f$ and $g$. The \emph{fundamental
group} $\pi_1(X,x_0)$ is the group of endpoint-fixed homotopy classes of
loops based at $x_0$, with concatenation.
\end{definition}

\begin{exercise}[Fundamental groups and winding]
Prove homotopy invariance of $\pi_1$ and the Seifert--van Kampen theorem for a
union with path-connected intersection. Compute
$\pi_1(S^1)\cong\Z$ by lifting loops to $\R$, and prove that the resulting
integer is the winding number.
\end{exercise}

\begin{definition}[Homology and cohomology]
The \emph{singular chain group} $C_n(X;R)$ is the free $R$-module on
continuous maps $\Delta^n\to X$, with alternating face boundary $\partial$.
The \emph{homology} is $H_n(X;R)=\ker\partial_n/\im\partial_{n+1}$. The
\emph{cochain complex} is $C^n(X;R)=\Hom(C_n(X;R),R)$ with coboundary
$\delta\phi=\phi\partial$, and its cohomology is $H^n(X;R)$.
\end{definition}

\begin{definition}[Products and pairings]
The \emph{Kronecker pairing} is
$H^n(X;R)\times H_n(X;R)\to R$, $([\phi],[c])\mapsto\phi(c)$. The
\emph{cup product} is the product on cohomology induced by the front--back
face formula on cochains.
\end{definition}

\begin{exercise}[Homology, cohomology, and degree]
Prove $\partial^2=0$, $\delta^2=0$, and homotopy invariance. Compute the
homology and cohomology rings of spheres and tori. For a continuous map
$f:S^n\to S^n$, define its degree by
$f_*[S^n]=\deg(f)[S^n]$ and prove multiplicativity and agreement with winding
number for $n=1$.
\end{exercise}

\begin{definition}[Fiber, vector, and principal bundles]
A \emph{fiber bundle} with fiber $F$ is a map $\pi:E\to X$ locally
isomorphic over $X$ to $U\times F\to U$. A \emph{vector bundle} has vector
spaces as fibers and linear transition maps. A \emph{principal $G$-bundle}
has a free right $G$-action and local equivariant trivializations. A
\emph{section} is a map $s:X\to E$ with $\pi s=\id_X$.
\end{definition}

\begin{exercise}[Bundles from transition functions]
Prove that transition functions satisfying the cocycle condition construct a
bundle and that changes of local trivialization give isomorphic bundles.
Prove that a principal bundle is trivial exactly when it has a global
section. Construct associated vector bundles from representations of $G$.
\end{exercise}

\begin{definition}[Global connections and holonomy]
A \emph{connection} on a vector bundle $E\to M$ is a map
$\nabla:\Gamma(E)\to\Omega^1(M;E)$ satisfying
$\nabla(fs)=df\otimes s+f\nabla s$. Its curvature is
$F_\nabla=\nabla^2\in\Omega^2(M;\End E)$. Its \emph{holonomy} along a loop is
the automorphism of the initial fiber given by parallel transport.
\end{definition}

\begin{exercise}[Curvature and holonomy]
Derive local connection and curvature forms and their gauge-transformation
laws. Prove the Bianchi identity. Show that infinitesimal holonomy is
curvature and compute the $U(1)$ holonomy around the boundary of a surface as
$\exp(i\int F)$ when a global potential is available.
\end{exercise}

\begin{definition}[Chern classes and character]
The \emph{Chern classes} of a complex vector bundle $E$ are the natural
classes $c_j(E)\in H^{2j}(X;\Z)$ determined by
$c(E\oplus F)=c(E)c(F)$ and the universal bundle. The \emph{Chern character}
is the natural ring map
$\operatorname{ch}:K^0(X)\to H^{\mathrm{even}}(X;\Q)$ whose value on a line
bundle $L$ is $e^{c_1(L)}$.
\end{definition}

\begin{exercise}[Curvature representatives]
For a Hermitian bundle with unitary connection, prove that
\[
\det\!\left(1+\frac{i}{2\pi}F_\nabla\right)
\quad\text{and}\quad
\Tr\exp\!\left(\frac{i}{2\pi}F_\nabla\right)
\]
are closed forms whose cohomology classes are independent of the connection.
Identify them with the Chern classes and Chern character and prove their
Whitney-sum formulas.
\end{exercise}

\begin{definition}[Berry geometry]
For a smooth family of one-dimensional spectral projections
$P_x$ on a Hilbert space, the ranges form a line bundle. The \emph{Berry
connection} is $\nabla=P\circ d$ and its curvature is the \emph{Berry
curvature}. Its holonomy is the \emph{Berry phase}. The integral
$\langle c_1(L),[\Sigma]\rangle$ is the \emph{Chern number} over an oriented
closed surface $\Sigma$.
\end{definition}

\begin{exercise}[Berry phase and quantized response]
Compute the Berry curvature of a two-level Hamiltonian
$H(n)=n\cdot\sigma$ over $S^2$ and its Chern number. For a gapped
two-dimensional periodic Hamiltonian, derive the Kubo formula expressing
Hall conductance as $e^2/h$ times the occupied-band Chern number. Deduce
integer quantization under gap-preserving perturbations and compare with
integer quantum Hall plateaus.
\end{exercise}

\section{K-Theory and Index}

\begin{definition}[Grothendieck completion and stable equivalence]
The \emph{Grothendieck group} of a commutative monoid $M$ is the universal
abelian group receiving a monoid map from $M$. Vector bundles $E,F$ are
\emph{stably equivalent} if $E\oplus\underline{\C}^m\cong
F\oplus\underline{\C}^n$ for some $m,n$. The group $K^0(X)$ is the
Grothendieck group of complex vector bundles on compact $X$; define
$K^1(X)=\widetilde K^0(\Sigma X)$.
\end{definition}

\begin{exercise}[Bott periodicity and topological phases]
Prove the clutching description of bundles on spheres. Establish Bott
periodicity in the form
$\widetilde K^0(S^{n+2})\cong\widetilde K^0(S^n)$ and compute the complex
$K$-groups of spheres and tori. Show that a gapped family of finite-rank
Hamiltonians defines the stable class of its occupied bundle and identify
the class-A invariants in dimensions one through three.
\end{exercise}

\begin{definition}[Operator K-theory]
For a unital $C^*$-algebra $A$, $K_0(A)$ is the Grothendieck group of stable
homotopy classes of projections in matrix algebras over $A$, and $K_1(A)$ is
the stable homotopy group of unitaries in matrix algebras over $A$.
\end{definition}

\begin{exercise}[Topological and operator K-theory]
Construct the Serre--Swan correspondence between vector bundles over compact
$X$ and finitely generated projective $C(X)$-modules. Deduce
$K^j(X)\cong K_j(C(X))$ and formulate Bott periodicity for $C^*$-algebras.
\end{exercise}

\begin{definition}[Compact and Fredholm operators]
An operator on a Hilbert space is \emph{compact} if it is a norm limit of
finite-rank operators. A bounded operator $T:H_0\to H_1$ is \emph{Fredholm}
if its kernel and cokernel are finite-dimensional and its range is closed;
its \emph{index} is
$\operatorname{ind}T=\dim\ker T-\dim\coker T$. The \emph{Calkin algebra} is
$\mathcal Q(H)=B(H)/\mathcal K(H)$.
\end{definition}

\begin{exercise}[Fredholm index and the Calkin algebra]
Prove that $T$ is Fredholm exactly when its image in the Calkin algebra is
invertible. Prove homotopy invariance, compact-perturbation invariance, and
additivity of the index. Compute the index of the unilateral shift and
identify the boundary map $K_1(\mathcal Q(H))\to K_0(\mathcal K(H))$.
\end{exercise}

\begin{definition}[Elliptic symbols]
For vector bundles $E,F$ over a compact smooth manifold, the \emph{principal
symbol} of an order-$m$ differential operator
$D:\Gamma(E)\to\Gamma(F)$ is the homogeneous bundle map
$\sigma_D:\pi^*E\to\pi^*F$ on $T^*M$. The operator is \emph{elliptic} if
$\sigma_D(x,\xi)$ is invertible for every $\xi\ne0$. Its \emph{analytic
index} is its Fredholm index on Sobolev completions; its \emph{topological
index} is the K-theoretic pushforward of its symbol class.
\end{definition}

\begin{exercise}[Atiyah--Singer index theorem]
Prove elliptic regularity and Fredholmness for elliptic operators on compact
manifolds. Establish the Atiyah--Singer equality
$\operatorname{ind}_{an}D=\operatorname{ind}_{top}(\sigma_D)$, using Bott
periodicity and the axiomatic characterization of the index. Derive the
Riemann--Roch and Dirac-index formulas as special cases.
\end{exercise}

\begin{definition}[Chirality and instantons]
Let $M$ be an oriented even-dimensional Riemannian spin manifold. The
\emph{chirality operator} splits its complex spinor bundle as
$S=S^+\oplus S^-$. For a Hermitian vector bundle $E$ with unitary connection
$A$, the coupled Dirac operator is odd and has a \emph{chiral part}
\[
D_A^+:\Gamma(S^+\otimes E)\longrightarrow\Gamma(S^-\otimes E).
\]
On an oriented Riemannian four-manifold, an $SU(r)$-connection is an
\emph{instanton} or \emph{anti-instanton} when $F_A=*F_A$ or $F_A=-*F_A$.
Its \emph{instanton number} is
\[
k=\int_M\operatorname{ch}_2(E)
  =-\frac{1}{8\pi^2}\int_M\Tr(F_A\wedge F_A)\in\Z.
\]
\end{definition}

\begin{exercise}[Fermion chirality detects instanton number]
Apply the Dirac-index formula from the preceding exercise to prove
\[
\operatorname{ind}D_A^+
=\dim\ker D_A^+-\dim\ker(D_A^+)^*
=\int_M\bigl[\widehat A(TM)\operatorname{ch}(E)\bigr]_4.
\]
For the fundamental $SU(2)$ bundle over $S^4$, show that this index is $k$
with the orientation convention in the definition. Use the Weitzenb\"ock
formula for a self-dual or anti-self-dual connection to determine which
chirality has no zero modes and deduce that the other has exactly $|k|$
zero modes. For a fermion in a representation $R$, derive the factor
$2T(R)k$. Finally, compute the fermionic measure under the axial rotation
$\psi\mapsto e^{i\alpha\gamma_5}\psi$ and obtain the integrated anomaly
$\Delta Q_5=2\,\operatorname{ind}D_A^+$, identifying the chirality change
carried by an instanton transition.
\end{exercise}

\section{Tensor Categories, Subfactors, and Topological Physics}

\begin{definition}[Monoidal and rigid categories]
A \emph{monoidal category} is a category $\mathcal C$ with a tensor functor,
unit object, and coherent associativity and unit isomorphisms. A \emph{dual}
of $X$ is an object $X^*$ with evaluation and coevaluation morphisms
satisfying the snake identities. A monoidal category is \emph{rigid} if
every object has left and right duals.
\end{definition}

\begin{definition}[Fusion, braided, and modular categories]
A \emph{fusion category} over $\C$ is a finite semisimple rigid
$\C$-linear monoidal category with simple unit. A \emph{braiding} is a
natural family $c_{X,Y}:X\otimes Y\to Y\otimes X$ satisfying the hexagon
identities. A ribbon fusion category is \emph{modular} if its braiding is
nondegenerate.
\end{definition}

\begin{exercise}[Fusion rules and braiding]
Prove that simple-object classes of a fusion category form a based ring under
tensor product. Derive the pentagon and hexagon equations for associativity
and braiding data. Compute quantum dimensions and braid-group
representations for the Fibonacci fusion rule
$\tau\otimes\tau\cong1\oplus\tau$.
\end{exercise}

\begin{definition}[Subfactors and conditional expectations]
A \emph{subfactor} is an inclusion $N\subseteq M$ of factors with common
unit. A \emph{conditional expectation} $E:M\to N$ is a positive unital
$N$-bimodule projection. For finite factors with normalized trace, the
\emph{Jones index} $[M:N]$ is the Murray--von Neumann dimension of
$L^2(M)$ as a left $N$-module.
\end{definition}

\begin{definition}[Basic construction]
On $L^2(M)$, the \emph{Jones projection} $e_N$ is the orthogonal projection
onto $L^2(N)$. The \emph{basic construction} is
$M_1=\langle M,e_N\rangle''$. Iteration gives the \emph{Jones tower}
$N\subseteq M\subseteq M_1\subseteq M_2\subseteq\cdots$.
\end{definition}

\begin{exercise}[Jones tower and Temperley--Lieb relations]
For finite-index $N\subseteq M$, derive
$e_Nxe_N=E(x)e_N$ and construct the canonical trace on the basic
construction. Prove that the normalized Jones projections in the tower
satisfy
\[
e_i^2=e_i,\qquad e_ie_j=e_je_i\ (|i-j|>1),\qquad
e_ie_{i\pm1}e_i=[M:N]^{-1}e_i.
\]
\end{exercise}

\begin{exercise}[Diagrammatic calculus and index quantization]
Construct the Temperley--Lieb diagrammatic category, with composition by
stacking and each closed loop valued at $\delta=[M:N]^{1/2}$. Use positivity
of its Markov trace to prove that index values below $4$ are
$4\cos^2(\pi/n)$ for integers $n\geq3$.
\end{exercise}

\begin{definition}[Braid groups and Markov traces]
The \emph{braid group} $B_n$ has generators $\sigma_i$ with relations
$\sigma_i\sigma_{i+1}\sigma_i=\sigma_{i+1}\sigma_i\sigma_{i+1}$ and
$\sigma_i\sigma_j=\sigma_j\sigma_i$ for $|i-j|>1$. A \emph{Markov trace} on
a tower of braid-group quotients is a compatible trace satisfying a fixed
stabilization rule.
\end{definition}

\begin{exercise}[Jones polynomial]
Construct a braid-group representation from the Temperley--Lieb generators.
Use the Markov trace and Markov moves to obtain an invariant of oriented
links, normalize it to the Jones polynomial, and derive its skein relation.
Compute it for the unknot, trefoil, and Hopf link.
\end{exercise}

\begin{definition}[Chern--Simons theory and Wilson loops]
For a connection $A$ on a principal bundle over an oriented
three-manifold and an invariant inner product on $\mathfrak g$, the
\emph{Chern--Simons functional} at level $k$ is
\[
S_{CS}(A)=\frac{k}{4\pi}\int_M
\left\langle A\wedge dA+\frac23A\wedge A\wedge A\right\rangle.
\]
For a representation $V$ and closed curve $K$, the \emph{Wilson loop} is
$W_V(K)=\Tr_V(\operatorname{Hol}_A(K))$.
\end{definition}

\begin{exercise}[Gauge invariance, level, and knot invariants]
Compute the change of $S_{CS}$ under a gauge transformation and derive the
integrality condition on $k$ required for $e^{iS_{CS}}$ to be gauge
invariant. Derive flatness as the field equation. Relate Wilson-loop
expectations for $SU(2)$ to the Jones polynomial at a root of unity.
\end{exercise}

\begin{definition}[Bordisms and topological field theories]
The \emph{bordism category} has closed oriented $(n-1)$-manifolds as objects,
oriented $n$-dimensional bordisms as morphisms, and disjoint union as tensor
product. An \emph{$n$-dimensional topological field theory} is a symmetric
monoidal functor from this category to vector spaces. An extended theory may
assign categories to codimension-two manifolds.
\end{definition}

\begin{exercise}[Functorial gluing]
Show that the monoidal-functor axioms give multiplicativity under disjoint
union, composition under gluing, duality under orientation reversal, and
mapping-class-group representations on state spaces. In two dimensions,
prove the equivalence with commutative Frobenius algebras.
\end{exercise}

\begin{definition}[Anyons and topological degeneracy]
In a two-dimensional gapped phase, an \emph{anyon type} is a simple object of
its braided fusion category. It is \emph{abelian} if its quantum dimension is
$1$ and \emph{nonabelian} otherwise. The \emph{topological degeneracy} on a
closed surface is the dimension of the state space assigned by the
associated topological field theory.
\end{definition}

\begin{exercise}[Abelian and nonabelian anyons]
Compute fusion spaces and braid actions for an abelian category and for the
Fibonacci category. Prove that nonabelian fusion spaces can have dimension
greater than one and that braiding acts projectively on them. Compute the
torus degeneracy from the simple anyon types in a modular category.
\end{exercise}

\begin{exercise}[Quantum Hall phases and measured topological order]
Use flux insertion for a Laughlin state at filling $\nu=1/m$ to derive
quasiparticle charge $e/m$, exchange phase $\pi/m$, Hall conductance
$\nu e^2/h$, and genus-$g$ degeneracy $m^g$. Compare these consequences
separately with conductance plateaus, fractional shot-noise charge
measurements, and interferometric phase data, stating which topological
predictions each observation tests and which it does not.
\end{exercise}

\end{document}
